\newif\iffull
\fulltrue
\iffull
\documentclass[a4paper,UKenglish,cleveref, autoref,
thm-restate,authorcolumns]{lipics-v2019}
\nolinenumbers
%This is a template for producing LIPIcs articles. 
%See lipics-manual.pdf for further information.
%for A4 paper format use option "a4paper", for US-letter use option "letterpaper"
%for british hyphenation rules use option "UKenglish", for american hyphenation rules use option "USenglish"
%for section-numbered lemmas etc., use "numberwithinsect"
%for enabling cleveref support, use "cleveref"
%for enabling autoref support, use "autoref"
%for anonymousing the authors (e.g. for double-blind review), add "anonymous"
%for enabling thm-restate support, use "thm-restate"

%\graphicspath{{./graphics/}}%helpful if your graphic files are in another directory
%\nolinenumbers
\bibliographystyle{plainurl}% the mandatory bibstyle

%\title{Maesuring Power of Star-Free Languages} Plevious title
\title{Measure-Theoretic Aspects of\\ Star-Free and Group Languages}

\titlerunning{Measure-Theoretic Aspects of Star-Free and Group Languages} %TODO optional, please use if title is longer than one line

\author{Ryoma Sin'ya}{Akita University, Akita, Japan}{ryoma@math.akita-u.ac.jp}{https://orcid.org/0000-0002-8152-998X}{JST ACT-X Grant Number JPMJAX210B.}
\author{Takao Yuyama}{Research Institute for Mathematical Sciences, Kyoto University, Kyoto, Japan}{yuyama@kurims.kyoto-u.ac.jp}{}{JSPS KAKENHI Grant Number JP20H05961.}
%TODO mandatory, please use full name; only 1 author per \author macro; first two parameters are mandatory, other parameters can be empty. Please provide at least the name of the affiliation and the country. The full address is optional

\authorrunning{R.\, Sin'ya} %TODO mandatory. First: Use abbreviated first/middle names. Second (only in severe cases): Use first author plus 'et al.'

\Copyright{Ryoma Sin'ya} %TODO mandatory, please use full first names. LIPIcs license is "CC-BY";  http://creativecommons.org/licenses/by/3.0/

\ccsdesc[500]{Theory of computation~Formal languages and automata theory}
%\ccsdesc[100]{\textcolor{red}{Replace ccsdesc macro with valid one}} %TODO mandatory: Please choose ACM 2012 classifications from https://dl.acm.org/ccs/ccs_flat.cfm 

\keywords{Automata, measure theory, star-free languages, group languages} %TODO mandatory; please add comma-separated list of keywords

\category{} %optional, e.g. invited paper

\relatedversion{} %optional, e.g. full version hosted on arXiv, HAL, or other respository/website
%\relatedversion{A full version of the paper is available at \url{...}.}

\supplement{}%optional, e.g. related research data, source code, ... hosted on a repository like zenodo, figshare, GitHub, ...

%\funding{(Optional) general funding statement \dots}%optional, to capture a funding statement, which applies to all authors. Please enter author specific funding statements as fifth argument of the \author macro.

%\acknowledgements{I want to thank \dots}%optional

%\nolinenumbers %uncomment to disable line numbering

%\hideLIPIcs  %uncomment to remove references to LIPIcs series (logo, DOI, ...), e.g. when preparing a pre-final version to be uploaded to arXiv or another public repository

%Editor-only macros:: begin (do not touch as author)%%%%%%%%%%%%%%%%%%%%%%%%%%%%%%%%%%
\EventEditors{John Q. Open and Joan R. Access}
\EventNoEds{2}
\EventLongTitle{42nd Conference on Very Important Topics (CVIT 2016)}
\EventShortTitle{CVIT 2016}
\EventAcronym{CVIT}
\EventYear{2016}
\EventDate{December 24--27, 2016}
\EventLocation{Little Whinging, United Kingdom}
\EventLogo{}
\SeriesVolume{42}
\ArticleNo{23}
%%%%%%%%%%%%%%%%%%%%%%%%%%%%%%%%%%%%%%%%%%%%%%%%%%%%%%

\usepackage{graphicx}
\usepackage{url}
\usepackage{ascmac,bm}
\usepackage{amsmath,amssymb}
\usepackage{float}
\usepackage{algorithm,algorithmic}
\usepackage{extarrows}
\usepackage{arydshln}
\usepackage{hyperref}
\usepackage{mleftright}
\usepackage{tikz, tikz-cd}
\usepackage{nicematrix}
\usepackage{mathtools}
\usepackage{centernot}

\theoremstyle{plain}
\newtheorem{conjecture}[theorem]{Conjecture}
\newtheorem{problem}[theorem]{Problem}
\else
\documentclass{llncs}

\title{Asymptotic Approximation\\ by Regular Languages}
\author{Ryoma Sin'ya\inst{1}}
\institute{Akita University\\
\email{ryoma@math.akita-u.ac.jp}\footnote{The author is also with RIKEN
AIP.}}

\usepackage[dvipdfmx]{graphicx}
\usepackage{url}
\usepackage{ascmac,bm}
\usepackage{amsmath,amssymb}
%	\theoremstyle{definition}
	\newtheorem{theo}{Theorem}
	\newtheorem{prop}{Proposition}
	\newtheorem{idea}{Idea}
	\newtheorem{lemm}{Lemma}
	\newtheorem{defi}{Definition}
	\newtheorem{ex}{Example}
\usepackage{float}
\usepackage{extarrows}
\usepackage{xcolor}

\bibliographystyle{splncs04.bst}
\fi

\newcommand{\CA}{{\cal A}}
\newcommand{\CJ}{\mathrel{\cal J}}
\newcommand{\CL}{\mathrel{\cal L}}
\newcommand{\CR}{\mathrel{\cal R}}
\newcommand{\CH}{\mathrel{\cal H}}

\newcommand{\ala}{{\it \`a la} }
\newcommand{\bala}{{\it \textbf{\`a la}} }

\newcommand{\pr}{\mathrm{Pr}}

\newcommand{\ie}{{\it i.e.}}
\newcommand{\eg}{{\it e.g.}}
\newcommand{\resp}{resp. }
\newcommand{\cf}{{\it cf. }}

\newcommand{\cesaro}{{Ces\`aro }}

\newcommand{\defeq}{=}
\newcommand{\defif}{\stackrel{\mathtt{def}}{\Longleftrightarrow}}

\newcommand{\zo}{\mathsf{ZO}}
\newcommand{\fin}{\mathsf{FIN}}
\newcommand{\dfin}{\mathsf{D}}
\newcommand{\rfin}{\mathsf{RD}}
\newcommand{\gfin}{\mathsf{GD}}
\newcommand{\group}{\mathsf{G}}
\newcommand{\gnil}{\mathsf{Gnil}}
\newcommand{\gsol}{\mathsf{Gsol}}
\newcommand{\gcom}{\mathsf{Gcom}}
\newcommand{\com}{\mathsf{Com}}
\newcommand{\modl}{\mathsf{MOD}}
\newcommand{\pt}{\mathsf{PT}}
\newcommand{\upol}{\mathsf{UPol}}
\renewcommand{\sf}{\mathsf{SF}}
\newcommand{\lt}{\mathsf{LT}}
\newcommand{\at}{\mathsf{AT}}
\newcommand{\gd}{\mathsf{GD}}
\renewcommand{\alph}{\mathrm{alph}}

\newcommand{\prim}{\mathsf{Q}}
\newcommand{\gold}{\mathsf{G}}
\newcommand{\dyck}{\mathsf{D}}
\newcommand{\pal}{\mathsf{P}}
\newcommand{\squ}{\mathsf{S}}
\newcommand{\kemp}{\mathsf{K}}
\newcommand{\othree}{\mathsf{O}_3}
\newcommand{\ofour}{\mathsf{O}_4}
\newcommand{\abcor}{\mathsf{L}}
\newcommand{\maj}{\mathsf{M}}

% algebra
\newcommand{\synt}{\mathsf{Synt}}
\newcommand{\ext}{\mathrm{Ext}}
\newcommand{\extreg}{\mathrm{Ext}(\reg)}
\newcommand{\rext}{\mathrm{RExt}}
\newcommand{\mext}{\mathrm{MExt}}

%group theory
\DeclareMathOperator{\UpTriMat}{UT}
\DeclareMathOperator{\WP}{WP}

%set theory
\newcommand{\symmdiff}{\bigtriangleup}
\newcommand{\comp}[1]{#1^{\mathrm{c}}}
\newcommand{\setmid}{\mathrel{}\middle|\mathrel{}}
%\newcommand{\setdelimiter}{:}
\newcommand{\setdelimiter}{\setmid}
\newcommand{\set}[2]{\mleft\{\, #1 \setdelimiter #2 \,\mright\}}

% density and measure
\newcommand{\pd}[1]{\delta_{#1}}
\renewcommand{\d}{\pd{A}}
\newcommand{\pref}{\mathtt{Pre}}
\newcommand{\suff}{\mathtt{Suf}}
\newcommand{\pfree}{\mathtt{pf}}
\newcommand{\sfree}{\mathtt{sf}}
\newcommand{\pcd}[1]{\delta_{#1}}
\newcommand{\cd}{\pcd{A}}
\renewcommand{\pm}[1]{\mu_{#1}}
\newcommand{\plm}[1]{\underline{\mu}_{#1}}
\newcommand{\pum}[1]{\overline{\mu}_{#1}}
\newcommand{\lm}{\plm{\reg}}
\newcommand{\um}{\pum{\reg}}
\newcommand{\m}{\pm{\reg}}

\newcommand{\rev}{\mathrm{rev}}

\renewcommand{\alph}{\mathtt{Alph}}

\usepackage{mathrsfs}
\newcommand{\bool}{{\mathscr B}}
\newcommand{\monc}{{\mathscr M}}
\newcommand{\nat}{\mathbb{N}}
\newcommand{\integer}{\mathbb{Z}}
\newcommand{\modint}[1]{\integer / #1 \integer}
\newcommand{\rat}{\mathbb{Q}}
\newcommand{\real}{\mathbb{R}}
\newcommand{\kk}{\mathbb{K}}
\newcommand{\card}[1]{\#\!\left(#1\right)}

\newcommand{\CC}{{\cal C}}
\newcommand{\CD}{{\cal D}}
\newcommand{\CB}{{\cal B}}
\newcommand{\CT}{{\cal T}}
\newcommand{\CO}{{\cal O}}

\newcommand{\unca}{\mathrm{UnCA}}
\newcommand{\ca}{\mathrm{CA}}
\newcommand{\reg}{\mathrm{REG}}
\newcommand{\cfl}{\mathrm{CFL}}
\newcommand{\dcfl}{\mathrm{DetCFL}}
\newcommand{\ucfl}{\mathrm{UnCFL}}

\newcommand{\slo}{<_{\text{lex}}}

\newcommand{\qedd}{\iffull
\else
\qed
\fi}



\begin{document}

\maketitle

\begin{abstract}
A language \(L\) is said to be \(\CC\)-measurable, where \(\CC\) is a
 class of languages, if there is an infinite sequence of languages in
 \(\CC\) that ``converges'' to \(L\).
In this paper, we investigate the measuring powers of \(\sf\) of the class of all star-free languages.
It is shown that the measuring power of star-free languages are coincides with the measuring power of generalised definite languages (a more restricted subclass of star-free languages), which implies the decidability of \(\sf\)-measurability for regular languages.
We also show that every regular language can be represented as a convergent sequence of Boolean combinations of star-free and group languages, \ie, every regular language is \(\bool(\sf \cup \group)\)-measurable where \(\group\) is the class of all group languages and \(\bool\CC\) is the Boolean algebra generated by \(\CC\).
We also demonstrate several consequences of the above two important facts.
\end{abstract}

%\keywords{automata, context-free languages, density, primitive words}
%\section{Prologue}
\section{Introduction}
One of the fundamental results in semigroup theory is the Krohn-Rhodes decomposition theorem stating that every finite monoid can be decomposed in a suitable sense into a wreath product of finite aperiodic monoids and finite groups.
According to the Krohn-Rhodes decomposition theorem, the class of \emph{star-free languages} (the language counterpart of aperiodic monoids) and \emph{group languages} (languages recognised by finite groups) can be considered as two major ``building blocks'' of regular languages.
While the theory of star-free languages is well-studied (\eg, Sch\"utzenberger's algebraic characterisation~\cite{Schutzenberger1965190},  McNaughton-Papert's logical characterisation~\cite{counterfree}, the dot-depth hierarchy (described later) proposed by Cohen and Brzozowski~\cite{dot} and recent progresses relating to this hierarchy (see \cite{dot3} or a survey \cite{dotdepth}), etc...), the theory of group languages seems not.
Although the definition of group languages is simple from an algebraic point of view, this class dose not have a language theoretic nor logical different characterisation.
To borrow the phrase used by Place--Zeitoun~\cite{place}: ``This makes it difficult to get an intuitive grasp about group languages, which may explain why this class remains poorly understood.''

This paper sheds new light on these two major building blocks --- star-free and group languages --- of regular languages, by using a measure-theoretic notion so-called \emph{\(\CC\)-measurability}.
A language \(L\) is said to be \(\CC\)-measurable if there exists an infinite sequence of pairs of languages \((K_n, M_n)_{n \in \nat}\) in \(\CC\) such that \(K_n \subseteq L \subseteq M_n\) holds for all \(n\) and the \emph{density} (described later) of the difference \(M_n \setminus K_n\) tends to zero as \(n\) tends to infinity.
The notion of \(\CC\)-measurability is introduced in \cite{S2021} and it was used for classifying non-regular languages by using regular languages.
This notion can be considered as a non-commutative and high-dimensional extension of Buck's measure theoretic notion for subsets of natural numbers so-called ``measure density''~\cite{Buck}, and \(\CC\)-measurability can be characterised by using a purely measure theoretic notion called \emph{Carath{\'e}odory extension}~\cite{carathe}.
% Approximating a complex object by more simple objects is a major concept
% in both computer science and mathematics.
% In the theory of formal languages, various types of approximations have
% been investigated (\eg,
% \cite{kappe2004,kappe2005,approx,CORDY2007125,CAMPEANU20013,Domaratzki}).
% For example, Kappes and Kintala~\cite{kappe2004} introduced 
% \emph{convergent-reliability} and \emph{slender-reliability} which
% measure how a given deterministic automaton \(\CA\) nicely approximates a given
% language \(L\) over an alphabet \(A\).
% A slightly modified version of approximation notion is
% \emph{bounded-\(\epsilon\)-approximation} which was introduced by Eisman and Ravikumar~\cite{Eisman2011}.
% They say that two languages \(L_1\) and \(L_2\) provide a
% bounded-\(\epsilon\)-approximation of language \(L\) if
% \(L_1 \subseteq L \subseteq L_2\) holds and
%  the ratio of their length-\(n\) difference satisfies
% \(
%  \card{(L_2 \setminus L_1) \cap A^n} / \card{A^n} \leq \epsilon
%  \)
% for every sufficiently large \(n \in \nat\).
% Perhaps surprisingly, they showed that no pair of regular languages can
%  provide a bounded-\(\epsilon\)-approximation of the
% language \(\{w \in \{a,b\}^* \mid w \text{ has more } a \text{'s than } b
% \text{'s}\}\) for any \(0 \leq \epsilon < 1\)~\cite{approx}.
% This result is a very strong \emph{in}approximable (by regular languages)
% example of certain non-regular languages.
% Also, there is a different framework of approximation so-called
% \emph{minimal-cover}~\cite{Domaratzki,CAMPEANU20013}, and
% a notion represents some \emph{in}approximability by regular languages so-called \emph{\(\reg\)-immunity}~\cite{immunity}.
% A model of approximation studied in this paper is \emph{\(\CC\)-measurability}, which was introduced in \cite{S2021}.
% This notion is rather close to bounded-\(\epsilon\)-approximation:
%  Instead of approximating by a \emph{single} regular language, we consider an approximation of some non-regular language \(L\) by an \emph{infinite sequence} of regular languages that ``converges'' to \(L\).
% While the membership problem for a given language \(L\) and a class
% \(\CC\) just asks whether \(L \in \CC\), the \(\CC\)-measurability asks the existence of
% an infinite sequence of languages in \(\CC\) that converges
% to \(L\).
% In this sense, measurability is much more difficult than the membership
% problem and its analysis is a challenging task.
So far, the decidability and alternative characterisation of \(\CC\)-measurability is systematically studied for some subclasses of regular languages.
For example, simple and decidable characterisations of \(\pt\)-measurable and \(\at\)-measurable languages are given in \cite{ltmeasure}, where \(\pt\) is the class of all piecewise testable languages and \(\at\) is the class of all alphabet testable languages (\ie, languages definable by the first-order logic with only one variable).
For the class \(\gd\) of \emph{generalised definite} languages (\ie, languages whose membership can be checked by finitely many prefix and suffix tests), an automata theoretic characterisation of the \(\gd\)-measurability for regular languages is given in \cite{gdmeasure} and it is decidable in quadratic time~\cite{ppl2024} with respect to the size of given deterministic automaton.
Some properties of the measuring power of star-free languages is investigated in \cite{carathe}, but the decidability for regular languages was left open.
The decidability of \(\group\)-measurability, where \(\group\) is the class of all group languages, for regular languages is still unknown, while the measurability of two restricted subclasses
 \(\gcom\) (commutative group languages) and \(\modl\) (languages defined by length-modulo condition) of group languages are shown to be decidable for regular languages~\cite{gcommeasure}. 

The main results of this paper (shown in Section~\ref{sec:sf}--\ref{sec:group}), relating to the measuring power of star-free and group languages, are three kinds as follows.
The first main result is about the measuring power of star-free languages (Section~\ref{sec:sf}).
It is shown that a language \(L\) is \(\sf\)-measurable if and only if \(L\) is \(\gd\)-measurable (Theorem~\ref{thm:sf}).
Since the \(\gd\)-measurability is known to be decidable for regular languages~\cite{gdmeasure}, this characterisation also gives the decidability of the \(\sf\)-measurability for regular languages.
We also gives a purely algebraic characterisation of the \(\sf\)-measurable regular languages (Theorem~\ref{thm:algsf}).
The second main result (Section~\ref{sec:ind}), which is an application of the first main result, is the \emph{probabilistic independence} of star-free and group languages (Theorem~\ref{thm:ind}) stating that
\(\cd(L \cap K) = \cd(L) \times \cd(K)\) holds for any
star-free language \(L\) and group language \(K\) where
\(\cd\) is the density function defined as
\(\cd(X) \defeq \lim_{n \rightarrow \infty} \frac{1}{n} \sum_{k = 0}^{n-1} \card{X \cap A^k}/\card{X^k}\) (where \(\card{S}\) denotes the cardinality of a set \(S\)).
Since the class of star-free languages corresponds to finite aperiodic monoids (monoids having no non-trivial subgroups) via Sch\"utzenberger's theorem~\cite{Schutzenberger1965190}, it is sometimes said that star-free languages and group languages are ``orthogonal.''
For example, it is described in \cite{nato} that ``This notion [aperiodic monoid] is in some sense ``orthogonal'' to the notion of groups'' and in \cite{DAM} that ``Aperiodic monoids are orthogonal to groups.'' Our second main result gives a rigorous statement how star-free and group languages are orthogonal in the probability theoretic sense.
Between \(\gd\) and \(\sf\), there is a fine-grained infinite strict hierarchy called the \emph{dot-depth hierarchy} introduced by Cohen and Brzozowski~\cite{dot} in 1971.
For a family \(\CC\) of languages, we denote by
\(\monc \CC = \{ L_1 \cdots L_k \mid k \geq 1, L_1, \ldots, L_k \in \CC
\} \cup \{\{\varepsilon\}\}\) the monoid closure of \(\CC\).
The dot-depth hierarchy starts with the family \(\CB_0\) of all finite
or co-finite languages, and continues as \(\CB_{i+1} = \bool \monc \CB_{i}\) for each \(i \geq 0\).
By definition, we have \(\sf = \bigcup_{i \geq 0} \CB_i\)
and also \(\CB_0 \subsetneq \gd \subsetneq \CB_1\).
However, the first main result tells us that this infinite strict hierarchy \emph{collapses} via taking all measurable sets in the sense that 
\(\{ L \mid L \text{ is } \gd\text{-measurable}\}
 = \{ L \mid L \text{ is } \sf\text{-measurable}\}\)
holds.
The last main result (Section~\ref{sec:group}) tells us that, for the measuring power on subclasses of group languages, the situation is exactly the opposite.
We show that, for any two local subvarieties\footnote{A \emph{local variety} is a family of regular languages of \(A^*\) closed under Boolean operations and quotients.} \(\CC \subsetneq \CD \subseteq \group\) of group languages, we have \(\{ L \mid L \text{ is } \CC\text{-measurable}\}
\subsetneq \{ L \mid L \text{ is } \CD\text{-measurable}\}\) (Theroem~\ref{thm:group-robust}).
%In Section~\ref{sec:fut} we summaries our results and explain some open problems.

\section{Preliminaries}\label{sec:pre}
This section provides the precise definitions of density and measurability.
We denote by \(\reg_A\) the family of all regular languages over an alphabet \(A\). For a word \(w \in A^*\) and a letter \(a \in A\), we write \(|w|_a\) for the number of occurrences of \(a\) in \(w\). The complement of \(L\) over \(A\) is denoted by \(\comp{L} = A^* \setminus L\).
We assume that the reader has a standard knowledge of algebraic language theory (\cf \cite{Lawson,DAM,MPRI}).

For a finite monoid \(M\) and an element \(x \in M\), the unique idempotent power of \(x\) is denoted \(x^\omega\).
A moniod \(M\) is called \emph{aperiodic} if \(x^{\omega+1} = x^\omega\) holds for any \(x \in M\).
Recall that Green's \(\CH\)-relation on a monoid \(M\) is defined as \(x \CH y \Leftrightarrow xM = yM \land Mx = My\), and
\(\CJ\)-relation is defined as \(x \CJ y \Leftrightarrow MxM = MyM\).
A monoid \(M\) is said to be \({\cal G}\)-trivial if every \({\cal G}\)-class is a singleton where \({\cal G} \in \{\CH, \CJ\}\).
It is well-known that a finite moniod \(M\) is aperiodic if and only if \(M\) is \(\CH\)-trivial.
The main targets of this paper are the following three subclasses of regular languages. The \emph{star-free languages}, the \emph{generalised definite languages} \(\gd\), and the \emph{group languages} \(\group\):
\begin{align*}
 \sf_A &\defeq \set{L \subseteq A^*}{L \text{ is recognised by a finite \(\CH\)-trivial monoid}},\\
 \gd_A &\defeq \bool\set{wA^*, A^*w}{w \in A^*},\\
 \group_A &\defeq \set{L \subseteq A^*}{L \text{ is recognised by a finite group}}.
\end{align*}
Also, we denote by \(\gcom_A\) the class of all languages recognised by finite commutative groups.
A monoid \(M\) divides a monoid \(N\) if \(M\) is a monoid homomorphic image of a submonoid of \(N\).
It is well-known that a monoid \(M\) recognises a language \(L\) if and only if the syntactic monoid \(M_L\) of \(L\) divides \(M\), hence
\(L\) is in \(\sf_A\) (\resp \(\group_A\)) if and only if \(M_L\) is a finite \(\CH\)-trivial monoid (\resp a finite group).
It is easy to see that a language \(L\) is generalised definite if and only if \(L = E \cup \bigcup_{(u, v) \in F} uA^*v\) for some finite sets \(E \subseteq A^*\) and \(F \subseteq A^* \times A^*\).
% \(L = E \cup \bigcup_{i \in I} F_i A^* G_i\) for some finite sets \(E\) and \(F_i, G_i \subseteq A^*\) for all \(i \in I\), where \(I\) is a finite index set.

\subsection{Density and measurability of formal language}
For a set $X$, we denote by $\card{X}$ the cardinality of $X$.
We denote by $\nat$ the set of natural numbers including $0$.

\begin{definition}[\cf \cite{codes}]\upshape
The \emph{density} \(\cd(L)\) of \(L \subseteq A^*\) is defined as
\[
 \cd(L) \defeq \lim_{n \rightarrow \infty} \frac{1}{n} \sum_{k = 0}^{n-1} \frac{\card{L \cap A^k}}{\card{A^k}}
\]
 if it exists, otherwise we write \(\cd(L) = \bot\).
 A language \(L\) is called \emph{null} if \(\cd(L) = 0\),
 and dually, \(L\) is called \emph{co-null} if \(\cd(L) = 1\).
\end{definition}

\begin{example}\label{ex:density}
It is known that every regular language has a rational density (\cf \cite{codes}) and it is computable.
For each word \(w\), the language \(A^* w A^*\),
 the set of all words that contain \(w\) as a factor,
 is of density one (\ie, co-null). This fact follows from the so-called {\em infinite monkey theorem}:
take any word \(w\). A random word of length \(n\) contains \(w\) as a factor with probability tending to one as \(n\) tends to infinity.
\end{example}

We list some basic properties of the density as follows.
\begin{lemma}[\cite{codes}]\label{lem:density}
 Let \(K, L \subseteq A^*\) with \(\cd(K) = \alpha, \cd(L) =
 \beta\). Then we have:
\begin{enumerate}
\item \(\cd(L \setminus K) = \beta - \alpha\) if \(K \subseteq L\). 
% \item \(\cd(\comp{K}) = 1 - \alpha\).
% \(\cd(K \cup L) \leq \alpha + \beta\) if \(\cd(K \cup L) \neq \bot\). 
\item \(\cd(K \cup L) = \alpha + \beta\) if \(K \cap L = \varnothing\). 
\item \(\cd(uLv) = \cd(L) \cdot \card{A}^{-|uv|}\) for each \(u, v \in A^*\).
\end{enumerate}
\end{lemma}
\begin{lemma}[{\cite[Example 13.4.9]{codes}}]\label{lem:groupdensity}
 Let \(\eta\colon A^* \rightarrow G\) be a surjective morphism onto a finite group. For any \(g \in G\), we have \(\cd(\eta^{-1}(g)) = \card{G}^{-1}\).
\end{lemma}

\begin{remark}\label{rem:kernel}
Let \(L\) be a regular language, \(M\) be a finite monoid and 
\(\eta: A^* \rightarrow M\) be a surjective morphism such that \(L = \eta^{-1}(\eta(L))\).
The \emph{kernel} \(K\) of \(M\) is a unique minimal ideal of \(M\) (it is always exists since \(M\) is finite),\ie, the minimal \(\CJ\)-class of \(M\).
For every \(m\) not in the kernel \(K\) of \(M\),
\(\eta^{-1}(m) \cap A^* w A^* = \emptyset\) holds for any \(w \in \eta^{-1}(K)\) hence \(\d(\eta^{-1}(m)) = 0\) by the infinite monkey theorem.
This implies that \(\d(\eta^{-1}(K)) = 1\).
Moreover, \(\d(\eta^{-1}(K)) = 1\) implies that there exist at least one element \(n \in K\) such that \(\d(\eta^{-1}(n)) > 0\) since \(K\) is finite.
Because \(K\) is the kernel of \(M\), for any \(m \in K\) there exists a pair of 
element \(x, y \in M\) such that \(xny = m\), hence \(\eta^{-1}(m)\) contains \(u \eta^{-1}(n) v\) where \(u \in \eta^{-1}(x)\) and \(v \in \eta^{-1}(y)\) so that
\(\d(\eta^{-1}(m)) \geq \d(\eta^{-1}(n)) \cdot \card{A}^{-(|u|+|v|)}\) by Lemma~\ref{lem:density}, \ie, \(\d(\eta^{-1}(m)) > 0\) holds for each \(m \in K\).
\end{remark}
For more detailed properties of the density function \(\cd\), see Chapter~13 of \cite{codes}.

The notion of ``measurability'' on formal languages is defined by a standard measure theoretic way as follows.
\begin{definition}[\cite{S2021}]\upshape
Let \(\CC_A\) be a family of languages over \(A\).
For a language \(L \subseteq A^*\), we define its
 \emph{\(\CC_A\)-inner-density} \(\plm{\CC_A}(L)\)
 and \emph{\(\CC_A\)-outer-density} \(\pum{\CC_A}(L)\)
 over \(A\) as
 \[
 \begin{alignedat}{2}
  \plm{\CC_A}(L) &\defeq& \, \sup &\set{\cd(K)}{K \subseteq L, K \in \CC_A, \cd(K) \neq \bot} \text{ and }\\
  \pum{\CC_A}(L) &\defeq& \inf &\set{\cd(K)}{L \subseteq K, K \in \CC_A, \cd(K) \neq \bot} \text{, respectively.}
 \end{alignedat}
 \]
A language \(L\) is said to be
 \emph{\(\CC_A\)-measurable} if
 \(\plm{\CC_A}(L) = \pum{\CC_A}(L)\) holds.
 We say that an infinite sequence \((L_n)_n\) of languages over \(A\)
 {\em converges to \(L\) from inner ({\em \resp}from outer)} if
 \(L_n \subseteq L\) (\resp \(L_n \supseteq L\)) for each \(n\) and \(\lim_{n \rightarrow \infty}\cd(L_n) = \cd(L)\).
A language \(L\) is \(\CC_A\)-measurable if and only if
there exist two convergent sequences of languages in \(\CC_A\) to \(L\) from inner and outer, respectively.
\end{definition}

The following is an example of \(\gcom\)-measurable non-regular languages.
For more detailed examples of \(\CC\)-measurable/immeasurable languages for a subclass \(\CC\) of regular languages, see \cite{S2021,carathe,gcommeasure}.

\begin{example}[\cite{S2021}]\label{ex:dyck2}
 The semi-Dyck language \(\dyck = \{\varepsilon, ab, aabb, abab, aaabbb, \ldots\}\) over \(A = \{a,b\}\) is \(\gcom\)-measurable.
 For each \(k \geq 2\), the language \(L_k = \{ w \in A^* \mid |w|_a = |w|_b \mod k \}\) is in \(\gcom\), because 
 \(\eta^{-1}(0) = L_k\) holds for the morphism \(\eta\colon A^* \rightarrow \modint{k}\) where \(\eta(a) = 1\) and \(\eta(b) = k-1\).
 Obviously, \(\dyck \subseteq L_k\) holds and it follows from Lemma~\ref{lem:groupdensity} that
 \(\cd(L_k) = 1/k\) holds. Hence \(\cd(L_k)\) tends to zero if \(k\) tends to infinity.
\end{example}

\subsection{Local varieties and profinite equations}
For a family \(\CC_A\) of languages over \(A\), we denote by
\(\ext(\CC_A)\) (\resp \(\rext(\CC_A)\)) the class of all \(\CC_A\)-measurable languages (\resp \(\CC_A\)-measurable \emph{regular} languages) over \(A\).
A family of regular languages over \(A\) is called a \emph{local variety} (\cf \cite{Adamek}) over \(A\) if it is closed under 
 Boolean operations and left-and-right quotients.
All classes \(\gd_A, \sf_A, \gcom_A\), and \(\group_A\) are local varieties.
The notion of \(\CC_A\)-measurability is well-behaved on Boolean operations and quotients as the following lemma says.
% A family of regular languages over \(A\)
%  is called a {\em local variety}~\cite{Adamek} over \(A\) if it is closed under 
%  Boolean operations and left-and-right quotients.
\begin{lemma}[\cite{carathe}]\label{lem:closure}
The operator \(\ext\) is a closure,
 \ie, it satisfies the following three properties for each \(\CC_A \subseteq \CD_A \subseteq 2^{A^*}\):
 (extensive)  \(\CC_A \subseteq \ext(\CC_A)\),
 (monotone) \(\ext(\CC_A) \subseteq \ext(\CD_A)\) and
 (idempotent) \(\ext(\ext(\CC_A)) = \ext(\CC_A)\).
Moreover, 
if \(\CC_A\) is a local variety, then \(\ext(\CC_A)\) is also a local variety.
\end{lemma}

Let \(\fin_A\) be the class of all finite and co-finite languages over \(A\).
Hereafter, we only consider an alphabet \(A\) such that \(\card{A} \geq 2\) (because \(\sf_A = \gd_A = \fin_A\) and \(\group_A = \gcom_A\) holds if \(\card{A} = 1\)) and sometimes omit the subscript \(A\) for denoting these classes of languages.

The notion of local varieties and local pseudovarieties can be captured by an appropriate ``equation'' notion called \emph{profinite equations} (\cf \cite{profinite}).
For a pair of words \(u\) and \(v\), we can define a dinstance between them by
\(d(u, v) = 2^{-r(u,v)}\)
where \(r(u,v) = \min\{ |\CA| \mid \CA \text{ is a deterministic automaton accepting one of } \{u, v\} \text{ but not the other } \}\).
Intuitively, \(r(u,v)\) represents the minimum number of states for separating \(u\) and \(v\) by using deterministic automata.
Then this function \(d\) is an ultrametric on \(A^*\), and the completion of the metric space \((A^*, d)\), denoted by \(\widehat{A^*}\), is called the set of profinite words over \(A\).
An element of \(\widehat{A^*}\) is a Cauchy sequence of (usual finite) words over \(A\) with respect to \(d\).
For example, the profinite word \(w^\omega\), associated with every word \(w\), can be defined as a limit of the sequence of words \((w^{n!})_{n \geq 1}\).
For any finite monoid \(M\) and morphism \(\eta: A^* \rightarrow M\), it is easy to see that \(\eta(w^{n!})\) is the idempotent power \(\eta(w)^{\omega}\) of \(\eta(w)\).
Hence, for a finite monoid \(M\), the equation \(w^{\omega+1} = w^{\omega}\) can be interpreted as \(x^{\omega+1} = x^\omega\) holds for any \(x \in M\), \ie, 
\(M\) is aperiodic (= \(\CH\)-trivial).
It is known that any local variety of languages can be captured by using profinite equations and vice versa.
We use the following theorem to prove Theorem~\ref{thm:group-robust} in Section~\ref{sec:group}.
%For the definition and basic properties of the free profinite monoid \(\widehat{A^*}\), see \cite{profinite,MPRI} for example.
\begin{theorem}[{\cite[Proposition 7.4]{Duality2008}}]\label{thm:duality}
  A set of regular languages over \(A\) is a Boolean quotienting algebra (\ie, local variety)
  if and only if it can be defined by a set of semigroup equations of the form \(u = v\), where \(u, v \in \widehat{A^*}\).
\end{theorem}
A regular language \(L \subseteq A^*\) \emph{satisfies} an equation \(u = v\), where \(u, v \in \widehat{A^*}\),
if the syntactic morphism \(\eta\colon A^* \to M_L\) of \(L\) satisfies \(\hat{\eta}(u) = \hat{\eta}(v)\),
where \(\hat{\eta}\colon \widehat{A^*} \to M_L\) is the unique uniformly continuous extension of \(\eta\)
(\cf \cite[Chapter X, Proposition 2.10]{MPRI}).

\begin{lemma}\label{lem:group-profinite-identity}
  For any equation \(u = v\) (where \(u, v \in \widehat{A^*}\)),
  there exists \(w \in \widehat{A^*}\) such that the equation \(w = \varepsilon\) is equivalent to \(u = v\) on any \(A\)-generated finite group.
\end{lemma}
\begin{proof}
  Let \((v_n)_{n \in \nat}\) be a Cauchy sequence in \(A^*\) converging to \(v\) in \(\widehat{A^*}\).
  The profinite word \(v^{\omega - 1} \in \widehat{A^*}\) is defined as \(v^{\omega - 1} = \lim_{n \to \infty} v_n^{n! - 1}\).
  Then, for every \(A\)-generated finite group \(\eta\colon A^* \to G\), we have \(\hat{\eta}(v^{\omega - 1}) = \hat{\eta}(v)^{-1}\) in \(G\).
  Thus \(w = u v^{\omega - 1}\) satisfies the condition of the lemma.
\end{proof}

\subsection{Local varieties of finite nilpotent/solvable groups}
This subsection provides background material for Section~\ref{sec:group}.
The reader may skip this subsection if interested only in the results up to Section~\ref{sec:ind}.

\begin{definition}
  For a positive integer \(n\), define
  \begin{align*}
    \gnil_A^{\leq n} &\defeq \set{L \subseteq A^*}{\text{\(L\) is recognised by a finite nilpotent group of class \(\leq n\)}},\\
    \gsol_A^{\leq n} &\defeq \set{L \subseteq A^*}{\text{\(L\) is recognised by a finite solvable group of derived length \(\leq n\)}}.
  \end{align*}
  Note that \(\gnil_A^{\leq 1} = \gsol_A^{\leq 1} = \gcom_A\).
\end{definition}
Recall that a group \(G\) is called a nilpotent group of class \(\leq n\) if and only if for any \(g_0, g_1, g_2, \dotsc, g_n \in G\),
the equation \([[\dotsm[[g_0, g_1], g_2]\dotsm], g_n] = 1_G\) holds (\eg, \cite[14.64]{hanna-neumann-variety}),
where \([g, h] \defeq g h g^{-1} h^{-1}\) is the commutator of \(g\) and \(h\).
Hence, for example, \(\gnil_A^{\leq 2}\) is a local variety defined by the set of equations of the form
\((u v u^{\omega - 1} v^{\omega - 1}) w (u v u^{\omega - 1} v^{\omega - 1})^{\omega - 1} w^{\omega - 1} = \varepsilon\) for \(u, v, w \in A^*\).
Similarly, each \(\gnil_A^{\leq n}\) is also a local variety.
It is also known that each \(\gsol_A^{\leq n}\) is defined by a set of equations (\eg, \cite[14.65]{hanna-neumann-variety}),
therefore \(\gsol_A^{\leq n}\) is also a local variety.

%For each \(n \geq 1\), it is known that there exists a finite nilpotent group of class \(n\) that is generated by at most \(n\) elements.
%For example, if \(p\) is a positive integer,
%then the group \(\UpTriMat(n, \integer / p \integer)\) of \(n \times n\) unitriangular matrices over the ring \(\integer / p \integer\)
%is a finite nilpotent group of class \(n - 1\), generated by \(n\) elements~\cite[Exercise 5.44.(ii)]{rotman-GTM148}.
%Therefore \(\gnil_A^{\leq n} \subsetneq \gnil_A^{\leq n + 1}\) for each \(n < \card{A}\).

%For each \(n \geq 1\), it is also known that there exists a finite solvable group of derived length \(n\)
%that is generated by at most \(n\) elements~\cite[9.22 Lemma]{rose-group}.
%For example, the \(n\)-fold iterated wreath product \((\integer / 2 \integer) \wr \dotsb \wr (\integer / 2 \integer)\)
%is a finite solvable group of derived length \(n\), generated by \(n\) elements.
%Therefore \(\gsol_A^{\leq n} \subsetneq \gsol_A^{\leq n + 1}\) for each \(n < \card{A}\).

%\begin{lemma}\label{lem:two-gen}
%  Any finite group \(G\) is generated by \(n\) elements as a group if and only if it is generated by \(n\) elements as a monoid.
%  In particular, if \(G\) is generated by two elements, there exists a surjective morphism \(\eta\colon A^* \to G\).
%\end{lemma}
%\begin{proof}
%  Since \(G\) is finite, the inverse of each \(g \in G\) is expressed as \(g^{\card{G} - 1}\).
%\end{proof}

The rest of this subsection is devoted to the proof of Proposition~\ref{prop:strict}.
\begin{proposition}\label{prop:strict}
  Suppose \(\card{A} \geq 2\).
  For every \(n \geq 1\), we have \(\gnil_A^{\leq n} \subsetneq \gnil_A^{\leq n + 1}\) and \(\gsol_A^{\leq n} \subsetneq \gsol_A^{\leq n + 1}\).
\end{proposition}
To prove Proposition~\ref{prop:strict}, we need Lemmata~\ref{lem:two-gen-nilpotent} and \ref{lem:two-gen-solvable} below.

\begin{lemma}\label{lem:two-gen-nilpotent}
  For \(n \geq 1\), there is a two-generated finite nilpotent group of class \(n\).
\end{lemma}
\begin{proof}
  For each \(n \geq 1\), the dihedral group \(D_{2^{n + 1}}\) of order \(2^{n + 1}\)
  is a two-generated finite nilpotent group of class \(n\) (see \eg, \cite[Chapter 4, \S 4]{suzuki-2-ja}).
\end{proof}

\begin{lemma}\label{lem:two-gen-solvable}
  For \(n \geq 1\), there is a two-generated finite solvable group of derived length \(n\).
\end{lemma}
To prove Lemma~\ref{lem:two-gen-solvable}, we use the following embedding theorem due to Neumann-Neumann,
and Lemmata~\ref{lem:solvable-length-n} and \ref{lem:solvable-reduce} below.

\begin{theorem}[{\cite[Theorem 5.6]{neumann-neumann-1959}}]\label{thm:embed-solvable}
  Every finite solvable group \(G\) of derived length \(n\) can be embedded in
  some two-generated finite solvable group \(H\) of derived length at most \(n + 2\).
\end{theorem}

%\begin{lemma}[{\cite[Lemma 2.1]{AM2011}}]\label{lem:length-wreath}
%  Let \(G\) be a solvable group of derived length \(n\) and \(k \geq 2\) be an integer.
%  Then the wreath product \(G \wr C_k\) is a solvable group of derived length \(n + 1\),
%  where \(C_k\) is the finite cyclic group of order \(k\).
%\end{lemma}

\begin{lemma}\label{lem:solvable-length-n}
  For each \(n \geq 1\), there exists a finite solvable group of derived length \(n\).
\end{lemma}
\begin{proof}
  %By Lemma~\ref{lem:length-wreath},
  %the \(n\)-fold iterated wreath product \((\integer / 2 \integer) \wr \dotsb \wr (\integer / 2 \integer)\)
  %is a finite solvable group of derived length \(n\).
  For example, the \(n\)-fold iterated wreath product \((\integer / 2 \integer) \wr \dotsb \wr (\integer / 2 \integer)\)
  is a finite solvable group of derived length \(n\)~\cite[9.22 Lemma]{rose-group}.
\end{proof}

\begin{lemma}\label{lem:solvable-reduce}
  If \(G\) is a solvable group of derived length \(n + 1\), then some factor group of \(G\) is solvable of derived length \(n\).
\end{lemma}
\begin{proof}
  Let \(G = G^{(0)} \rhd G^{(1)} \rhd \dotsb \rhd G^{(n)} \rhd G^{(n + 1)} = \{1_G\}\) be the derived series of \(G\),
  where \(G^{(i + 1)} \defeq [G^{(i)}, G^{(i)}]\).
  Since each \(G^{(i + 1)}\) is a characteristic subgroup of \(G^{(i)}\), the group \(G^{(n)}\) is also a characteristic subgroup of \(G\),
  hence in particular a normal subgroup of each \(G^{(i)}\).
  Let \(\bar{G} \defeq G / G^{(n)}\).
  Then \(G^{(0)} / G^{(n)} \rhd G^{(1)} / G^{(n)} \rhd \dotsb \rhd G^{(n)} / G^{(n)} = \{1\}\) is the derived series of \(\bar{G}\) of length \(n\),
  since we have an isomorphism \([G^{(i)} / G^{(n)}, G^{(i)} / G^{(n)}] \cong [G^{(i)}, G^{(i)}] / G^{(n)} = G^{(i + 1)} / G^{(n)}\).
\end{proof}

\begin{proof}[Proof of Lemma~\ref{lem:two-gen-solvable}]
  By Lemma~\ref{lem:solvable-length-n}, there exists a finite solvable group \(G\) of derived length \(n\).
  By Theorem~\ref{thm:embed-solvable},
  \(G\) can be embedded in some two-generated finite group \(H\) of derived length at most \(n + 2\).
  %By Lemma~\ref{lem:length-wreath}, the derived length of \(G \wr C_k\) is \(n + 1\).
  In particular, the derived length of \(H\) is one of \(n\), \(n + 1\), or \(n + 2\).
  Thus, by Lemma~\ref{lem:solvable-reduce},
  we obtain a two-generated finite solvable group of derived length \(n\) as a factor group of \(H\).
\end{proof}

Lemma~\ref{lem:WP-synt} below is the last piece of the proof of Proposition~\ref{prop:strict}.
%Let \(G\) be a group and \(\eta\colon A^* \to G\) be a surjective morphism.
%The \emph{word problem} \(\WP_\eta(G)\) of \(G\) (with respect to \(\eta\))
%is the language over \(A\) defined as \(\WP_\eta(G) \defeq \eta^{-1}(1_G)\).
\begin{lemma}[{\cite[3.4.9 Proposition]{Holt2017}}]\label{lem:WP-synt}
  Let \(G\) be a group and \(\eta\colon A^* \to G\) be a surjective morphism.
  Then the syntactic monoid of the language \(\eta^{-1}(1_G) \subseteq A^*\) is isomorphic to \(G\).
  %Then the syntactic monoid of \(\WP_\eta(G)\) is isomorphic to \(G\).
  %In particular, every \(A\)-generated group is the syntactic monoid of some language over \(A\).
\end{lemma}
The language \(\eta^{-1}(1_G)\) is known as the \emph{word problem} for the group \(G\) (over \(A\)).

\begin{proof}[Proof of Proposition~\ref{prop:strict}]
  By Lemma~\ref{lem:two-gen-nilpotent}, there exists a two-generated finite nilpotent group \(G\) of class \(n + 1\).
  Since \(G\) is finite and generated by two elements as a group, it is generated by the two elements as a monoid.
  Hence there exists a surjective morphism \(\eta\colon A^* \to G\).
  Since \(\eta^{-1}(1_G)\) is clearly recognised by \(G\), we have \(\eta^{-1}(1_G) \in \gnil_A^{\leq n + 1}\).
  Contrary, suppose \(\eta^{-1}(1_G) \in \gnil_A^{\leq n}\).
  Then \(\eta^{-1}(1_G)\) is recognised by some finite nilpotent group \(H\) of class \(n\),
  hence the syntactic monoid of \(\eta^{-1}(1_G)\), which is isomorphic to \(G\) by Lemma~\ref{lem:WP-synt}, divides \(H\).
  In general, if \(K\) is a nilpotent group of class \(c\),
  then every subgroup of \(K\) and every factor group of \(K\) is nilpotent of class at most \(c\).
  Hence \(G\) is nilpotent of class at most \(n\), contradiction.
  (Note that every submonoid of a group is a subgroup.)
  Thus \(\gnil_A^{\leq n} \subsetneq \gnil_A^{\leq n + 1}\).

  By Lemma~\ref{lem:two-gen-solvable}, there exists a two-generated finite solvable group \(G\) of derived length \(n + 1\).
  Then one can similarly show \(\eta^{-1}(1_G) \in \gsol_A^{\leq n + 1} \setminus \gsol_A^{\leq n}\) for a surjective morphism \(\eta\colon A^* \to G\).
\end{proof}


\section{Measuring Power of Star-Free Languages}\label{sec:sf}

In this section, we give a proof of the first main theorem, which gives a decidable characterisation of \(\sf\)-measurability for regular languages.
The key lemma of the proof of Theorem~\ref{thm:sf} is the following lemma.
\begin{lemma}\label{lem:cat}
 The class \(\ext(\gd)\) is closed under concatenation.
\end{lemma}

\begin{theorem}\label{thm:sf}
The class of star-free languages and the class of generalised definite languages have the same measuring power, \ie, \(\ext(\sf) = \ext(\gd)\).
\end{theorem}
\begin{proof}[Proof of Theorem~\ref{thm:sf}]
By definition, \(\fin\) is contained in \(\gd\).
Moreover, \(\ext(\gd)\) is closed under Boolean operations by Lemma~\ref{lem:closure} since \(\gd\) is so.
Combining these facts with Lemma~\ref{lem:cat}, we can conclude that
\(\sf \subseteq \ext(\gd)\), which implies 
\(\ext(\sf) \subseteq \ext(\ext(\gd)) = \ext(\gd) \subseteq \ext(\sf)\)
by Lemma~\ref{lem:closure}.
That is, \(\ext(\sf) = \ext(\gd)\) holds.
\end{proof}
\begin{corollary}
 The \(\sf\)-measurability for a given deterministic automaton is decidable in quadratic time.
\end{corollary}

The proof of Lemma~\ref{lem:cat} requires some additional notion and lemma.
For a given generalised definite language \(L = F A^* G\) where \(F, G\) are finite, we define its \emph{prefix density}
\(\pred(L)\) as \(\d(F A^*)\) and its
\emph{suffix density} \(\sufd(L)\) as \(\d(A^* G)\).
By Lemma~\ref{lem:density}, we have \(\d(L) = \pred(L) \cdot \sufd(L)\).
%\section{Measure-Theoretic Decomposition Theorem of Regular Languages}

\begin{theorem}\label{thm:dec}
Every regular language is \(\bool(\sf \cup \group))\)-measurable, \ie, \(\reg = \rext(\bool(\sf \cup \group))\).
\end{theorem}

\begin{corollary}
A language \(L\) is \(\reg\)-measurable if and only if 
\(L\) is \(\bool(\sf \cup \group))\)-measurable.
\end{corollary}
\section{Some Applications of Two Main Theorems}
\subsection{Axiomatic characterisation of the density function \(\d\)}
\subsection{Probabilistic independence of star-free and group languages}
\subsection{A simple proof of the existence of non-primitive words in a regular language with positive density}
\section{Hierarchies of Group Languages Are Not Collapsed by \(\rext\)}\label{sec:group}

The main result of this section is Corollary~\ref{cor:strict-nilp-solv},
which is a consequence of Theorem~\ref{thm:group-robust} below and states that
the classes of regular languages approximable by nilpotent/solvable group languages form strict hierarchies.
This is completely different from the situation in Theorem~\ref{thm:sf},
which shows that the dot-depth hierarchy is collapsed by \(\rext\).

\begin{theorem}[Y.]\label{thm:group-robust}
  If \(\CC \subseteq \group_A\) is a local variety of group languages, then \(\rext(\CC) \cap \group_A = \CC\).
  In particular, if \(\CC \subsetneq \CD \subseteq \group_A\) are local subvarieties,
  then \(\rext(\CC) \cap \CD = \CC\), \(\CD \not\subseteq \rext(\CC)\), and \(\rext(\CC) \subsetneq \rext(\CD)\).
\end{theorem}
Our previous result \cite[Theorem 1]{gcommeasure} is the special case of Theorem~\ref{thm:group-robust} where \(\CC = \gcom_A\).

\begin{proof}
  Let \(L \in \group_A \setminus \CC\).
  It suffices to show that \(L\) is \(\CC\)-immeasurable.
  Since \(\CC\) is a local variety and \(L \notin \CC\), Theorem~\ref{thm:duality} implies
  that there exists an equation \(u = v\) (where \(u, v \in \widehat{A^*}\)) that every language in \(\CC\) satisfies but \(L\) does not.
  Let \(G\) be the syntactic monoid of \(L\), which is a group, and \(\eta\colon A^* \to G\) be the syntactic morphism of \(L\).
  By Lemma~\ref{lem:group-profinite-identity},
  there exists \(w \in \widehat{A^*}\) such that the equation \(w = \varepsilon\) is equivalent to \(u = v\) on any \(A\)-generated finite group.
  Hence, for every \(L_1 \in \CC\) and its syntactic morphism \(\theta_1\colon A^* \to G_1\),
  we have \(\hat{\theta}_1(w) = \hat{\theta}_1(\varepsilon) = 1_{G_1}\) and \(\hat{\eta}(w) \neq \hat{\eta}(\varepsilon) = 1_G\).
  Since \(\hat{\eta}(w) \neq 1_G\) and \(G\) is the syntactic monoid of \(L\), there exist \(x, y \in A^*\)
  such that
  \begin{equation}
    \eta(x) \hat{\eta}(w) \eta(y) \in \eta(L) \centernot\iff \eta(x y) \in \eta(L). \label{eq:ebc45b27}
  \end{equation}
  Since the two languages \(\eta^{-1}(\eta(x)), \eta^{-1}(\eta(y)) \subseteq A^*\) are group languages,
  they are non-null by Lemma~\ref{lem:groupdensity}.
  Hence the concatenation \(I \defeq \eta^{-1}(\eta(x)) \eta^{-1}(\eta(y))\) is also non-null (by Lemma~\ref{lem:density}, for example).
  To complete the proof, it suffices to show that \(L_1, L_2 \in \CC\) and \(L_1 \subseteq L \subseteq L_2\) implies \(I \subseteq L_2 \setminus L_1\).
  Let \(\theta_1\colon A^* \to G_1\) and \(\theta_2\colon A^* \to G_2\) be the syntactic morphism of \(L_1\) and \(L_2\), respectively.
  Let \((w_n)_{n \in \nat}\) be a Cauchy sequence in \(A^*\) converging to \(w\) in \(\widehat{A^*}\).
  Then we have \(\eta(w_k) = \hat{\eta}(w)\), \(\theta_1(w_k) = \hat{\theta}_1(w)\), and \(\theta_2(w_k) = \hat{\theta}_2(w)\) for sufficiently large \(k \in \nat\).

  We show \(I \subseteq L_2 \setminus L_1\).
  Let \(z \in I\).
  Then there exist \(\tilde{x} \in \eta^{-1}(\eta(x))\) and \(\tilde{y} \in \eta^{-1}(\eta(y))\) such that \(z = \tilde{x} \tilde{y}\).
  Let \(i \in \{1, 2\}\).
  Since \(\theta_i(w_k) = 1_{G_i}\) and \(G_i\) is the syntactic monoid of \(L_i\),
  we have \(\theta_i(\tilde{x} w_k \tilde{y}) \in \theta_i(L_i) \iff \theta_i(\tilde{x} \tilde{y}) \in \theta_i(L_i)\),
  \ie,
  \begin{equation}
    z = \tilde{x} \tilde{y} \in L_i \iff \tilde{x} w_k \tilde{y} \in L_i. \label{eq:fbaddcb1}
  \end{equation}
  It follows from \eqref{eq:ebc45b27}, \(\eta(\tilde{x}) = \eta(x)\), \(\eta(\tilde{y}) = \eta(y)\), and \(\hat{\eta}(w) = \eta(w_k)\) that
  \begin{equation}
    z = \tilde{x} \tilde{y} \in L \centernot\iff \tilde{x} w_k \tilde{y} \in L. \label{eq:500d85c6}
  \end{equation}
  %It is easy to show that \eqref{eq:fbaddcb1}, \eqref{eq:500d85c6},
  %and \(L_1 \subseteq L \subseteq L_2\) implies \(z \in L_2 \setminus L_1\).
  We show \(z \in L_2 \setminus L_1\) by considering two cases depending on whether \(z \in L\) or \(z \notin L\).
  First, suppose that \(z \in L\).
  Since \(L \subseteq L_2\), we have \(z \in L_2\).
  On the other hand, by \eqref{eq:500d85c6}, we know that \(\tilde{x} w_k \tilde{y} \notin L \supseteq L_1\), which implies \(\tilde{x} w_k \tilde{y} \notin L_1\).
  Then, by \eqref{eq:fbaddcb1}, it follows that \(z \notin L_1\), and thus we obtain \(z \in L_2 \setminus L_1\).
  The case where \(z \notin L\) can be proved in a similar manner.
\end{proof}

\begin{corollary}\label{cor:strict-nilp-solv}
  Suppose \(\card{A} \geq 2\).
  For \(n \geq 1\), we have
  \begin{itemize}
    \item \(\rext(\gnil_A^{\leq n}) \subsetneq \rext(\gnil_A^{\leq n + 1}) \subsetneq \rext(\group_A)\), and
    \item \(\rext(\gsol_A^{\leq n}) \subsetneq \rext(\gsol_A^{\leq n + 1}) \subsetneq \rext(\group_A)\).
  \end{itemize}
\end{corollary}
\begin{proof}
  By Proposition~\ref{prop:strict} and Theorem~\ref{thm:group-robust}.
\end{proof}


\section{Conclusion and Open Problems}\label{sec:fut}
Theorem~\ref{thm:sf} and our previous work~\cite{gdmeasure,ppl2024} give a decidable characterisation of the \(\sf\)-measurable regular languages.
Its algebraic characterisation (Theorem~\ref{thm:algsf}) can be considered as a natural extension of Sch\"utzenberger's theorem for star-free languages.
The decidability of the \(\group\)-measurability for regular languages is still open.
Perhaps the profinite method can be useful to solve this problem.
The probabilistic independence of star-free and group languages describes how these two classes are orthogonal.
%It is interesting that whether there exists a local variety of regular languages \(\CC\) properly containing \(\sf\) and the probabilistic independence of \(\CC\) and \(\group\) also holds.
%If there is no such class, one can regard that the class of star-free languages is a ``maximal'' subclass of regular languages which is orthogonal to the class of group languages in the probabilistic sense.

In 2023, the first author posed a question whether the equation
\(\reg = \bool(\rext(\sf) \cup \rext(\group))\) holds or not, at some domestic workshops in Japan.
If this kind of equation \emph{were} true, it can be considered as a ``measure-theoretic decomposition theorem'' of regular languages.
Later, the second author claimed that a weaker equation
\(\reg = \rext(\bool(\sf \cup \group))\) holds and
this result was announced at the international conferences Highlights 2024 by the first author.
However, after some discussion between these two authors, a flaw was found in the proof of \(\reg = \rext(\bool(\sf \cup \group))\).
Actually, this equation does \emph{not} hold as follows.
\begin{proposition}\label{prop:counterexample}
Let \(\CT_4\) be the full transformation monoid on the four elements set \(\{1,2,3,4\}\) whose identity element is the identity function \(\mathrm{id}\).
Define \(e, f \in \CT_4\) as
\begin{align*}
 e\colon 1 \mapsto 1, 2 \mapsto 2, 3 \mapsto 1, 4 \mapsto 2 \quad \text{ and } \quad
 f\colon 1 \mapsto 4, 2 \mapsto 3, 3 \mapsto 3, 4 \mapsto 4.
\end{align*}
Let \(M\) be the submonoid of \(\CT_4\) generated by \(e\) and \(f\).
For an alphabet \(A = \{a,b\}\), define a morphism \(\eta\colon A^* \rightarrow M\) as \(\eta(a) = e\) and \(\eta(b) = f\).
Then \(\eta^{-1}(e)\) is \(\bool(\sf_A \cup \group_A)\)-immeasurable.
In particular, we have \(\reg_A \supsetneq \rext_A(\bool(\sf_A \cup \group_A))\).
\end{proposition}
The authors are still interested in whether there is a suitable notion \(\CO\colon 2^{A^*} \rightarrow 2^{A^*}\) of ``operation'' on languages, which is richer than the Boolean closure operator \(\bool\) but the simpler than the operation on languages corresponds to the wreath product operation on monoids, for a measure-theoretic decomposition theorem of regular languages like
 \(\reg = \CO(\rext(\sf) \cup \rext(\group))\) or
 \(\reg = \rext(\CO(\sf \cup \group))\).


\bibliography{ref}

\end{document}
